\chapter{Threats to Validity} \label{chap:threats}

\todo{Identifies threats to validity of this research. These are the kinds of threats you should discuss:
1. Internal Validity: Concerns whether the research design and conduct accurately capture what the study aims to investigate. This includes issues like confounding variables, biases in data collection or analysis, and errors in experimental design.
2. External Validity: Refers to the generalisability of the findings. It questions whether the results can be applied to contexts outside of the study, such as different software systems, user groups, or environments.
3. Construct Validity: Deals with the appropriateness of the measurements and procedures used in the study. It examines whether the study actually measures what it intends to measure.
4. Conclusion Validity: Relates to the degree to which conclusions drawn from the study are credible or believable. This involves the statistical or logical strength of the study's findings.
Addressing these threats involves not only identifying and discussing them but also, where possible, explaining how the research design minimizes their impact. This demonstrates the rigor and thoroughness of the research and builds trust in the study's conclusions.
Not all are applicable in all papers, but we expect the student to think about what threats are applicable to the particular research. About one page.}
\Blindtext[3]
